\chapter{Implementazione del Framework}
In questo capitolo viene esposto l'intero framework di compressione proposto, costruito sulle fondamenta teoriche esposte nel Capitolo \ref{cap:teoria}.

\section{Gestione e Divisione dei dati}
Nella fase iniziale il framework di compressione, prevede di istanziare il dataset. In particolare, inizialmente viene eseguita la suddivisione in \textit{Dati di Train} e \textit{Dati di Test}, successivamente i primi vengono ulteriormente divisi in \textit{Training Set} e \textit{Validation Set}. Il secondo verrà invece utilizzato esclusivamente per monitorare l'accuratezza del modello durante le iterazioni di compressione. Per quanto concerne il \textit{Training Set}, questo viene utilizzato alla fine di ogni iterazione per eseguire una serie di epoche di \textit{Fine Tuning}, al fine di permettere un riadattamento dei parametri del modello, ma anche come sorgente da cui estrarre il \textit{Search Set}. \newline
Quest'ultimo è formato da un numero ben preciso di campioni per classe $N$, che costituisce un iper-parametro del processo di compressione, in questo modo la sua dimensione corrisponde ad $N \times |Classes|$. Dagli esperimenti effettuati emerge un valore ottimale di $N$ di 25 campioni per classe, poiché permette una buona approssimazione dei gradienti, ma allo stesso tempo consente di mantenere contenuto il tempo necessario alla fase di ricerca \textit{NAS}. Infatti, il \textit{Search Set}, viene utilizzato su ogni nodo dell'albero di ricerca (ovvero su ogni \textit{Subnet}) sia per il calcolo della metrica di importanza per ogni parametro (che prevede quindi un passo di \textit{backpropagation}, dal grande costo computazionale), sia per la valutazione dell'accuratezza. \newline 
Per evitare che la fase di ricerca, utilizzando un unico \textit{Search Set} per tutte le iterazioni, provochi un \textit{\textbf{Overfitting Architetturale}}, con conseguenti difficoltà di generalizzazione del modello compresso, il \textit{Search Set}, viene \textbf{campionato ad ogni iterazione} a partire dal \textit{Training Set}, fornendo campioni differenti all'algoritmo \textit{NAS}, di iterazione in iterazione. Questo permette di ottenere un modello risultante che è in grado di generalizzare al meglio sul problema di classificazione, proprio perché l'architettura finale, è il risultato ottenuto con una maggiore quantità e varietà di dati. \newline
Un altro contributo fondamentale al mantenimento della capacità di generalizzazione del modello, è dato dalla \textbf{Data Augmentation}, che viene sempre applicata al \textit{Search Set}, ed è la stessa che viene applicata ai dati del \textit{Training Set}. Si potrebbe pensare che questo introduca rumore nei gradienti, e quindi influenzi la precisione e l'affidabilità della metrica di importanza. Nella pratica però il problema del rumore è fortemente limitato dall'utilizzo della media dei gradienti sui batch, per il calcolo della metrica di importanza, com'è possibile osservare nella formula \ref{eq:metrica_importanza}. La \textit{Data Augmentation} applicata al \textit{Search Set} è risultata molto efficace, nei test effettuati, per impedire che venissero sacrificate delle dimensioni del modello essenziali per mantenere la sua capacità di generalizzazione, traducendosi in questo modo in valori di accuratezza sul \textit{Test Set} e sul \textit{Validation Set} più elevati durante tutta l'esecuzione del framework. \newline
Notare che è possibile interpretare il framework di compressione come un \textbf{ottimizzazione a due livelli} del modello. In particolare il livello superiore è occupato dal processo di ricerca \textit{NAS}, che esegue un vero e proprio \textbf{training architetturale}, mentre il livello inferiore è quello del \textit{Fine Tuning}, che invece si occupa di \textbf{ottimizzare i parametri}, date le modifiche architetturali operate dal livello superiore. Questa visione del metodo permette di comprendere al meglio le motivazioni dietro la specifica gestione dei dati, ed in particolare del \textit{Search Set}, infatti la dinamicità e l'\textit{Augmentation} di quest'ultimo sono fattori che contribuiscono a scongiurare il fenomeno dell'\textit{overfitting architetturale}, esattamente come tecniche quali l'\textit{early stopping}, i \textit{Learning Rate Scheduler}, ecc\dots vengono impiegati per prevenire l'\textit{overfitting} dei parametri del modello.

\section{Framework di Compressione Iterativo}
Una volta definita la strategia di gestione dei dati, è necessario definire la struttura del framework di compressione iterativo, in particolare cosa avviene ad ogni iterazione. \newline
Inizialmente viene campionato il \textit{Search Set}, in modo tale da poterlo utilizzare in fase di ricerca \textit{NAS}. Successivamente inizia la prima fase dell'iterazione: il \textbf{Pruning}. La prima operazione che viene eseguita è l'algoritmo di ricerca \textit{BnB}, che verrà analizzato nella sezione \ref{section:metodo_ricerca}, il quale restituisce lo stato migliore individuato, ovvero una maschera di \textit{pruning} per l'intero modello. \newline
A questo punto la maschera viene applicata al modello, e questo viene successivamente ricostruito istanziando tutti i tensori dei parametri ed effettuando una copia solo di quelli non tagliati in fase di ricerca. Viene quindi creato un nuovo modello \textit{ViT}, del quale sono state effettivamente ridotte le dimensioni dei tensori dei parametri e questo risulta fondamentale perché implica che il framework proposto non si limita ad applicare delle maschere di \textit{pruning}, che non avrebbero alcun impatto sulle prestazioni effettive del modello sull'hardware, bensì produce dei nuovi modelli di dimensioni realmente minori, con prestazioni realmente superiori. In particolare viene ridotta la latenza in fase di inferenza ma anche la memoria necessaria (VRAM) per utilizzare il modello.

\section{Algoritmo di Ricerca adottato}
\label{section:metodo_ricerca}